% \iffalse meta-comment
%
% hit-fonts.dtx 是 book / art / artplus 共用的字体层
%
% \fi
%
% \section{共享字体层}
%
% \changes{v3.2a}{0000/00/00}{字体层从三份各自的 \texttt{*-fonts} 模块抽出合并}
% book、art、artplus 的字体设置原来是三份拷贝，只差 \cs{normalsize} 的行距。
% 这里合成一份，\cs{normalsize} 仍旧留在各自的类里。
%
% 这个文件必须写在每个 \cs{file} 列表的第二位，也就是类文件头之后、类自己的 dtx
% 之前。因为 docstrip 按输入文件首次出现的顺序排全局次序，后出现的排在后面，写在
% 别处会报 Incompatible order of input files。
%
% 第二位在 \cs{LoadClass} 之前，此时 \pkg{fontspec} 和 \pkg{xeCJK} 还没加载，
% \option{font} 与 \option{cjk-font} 也还没解析。所以这个文件只做声明和定义，要
% 执行的部分放进 \cs{hit@setup@fonts}，由各个类在原来 \texttt{*-fonts} 的位置调用。
%
%    \begin{macrocode}
%<*shared-fonts>
\bool_new:N \g__hit_platform_mac_bool
\str_new:N \l__hit_font_str
\str_new:N \l__hit_cjk_font_str
\cs_new_protected:Npn \hit@setup@fonts
  {
%    \end{macrocode}
%
% 平台探测。\pkg{expl3} 只能判断是不是 Windows，macOS 看系统字体文件在不在。
%    \begin{macrocode}
\sys_if_platform_unix:T
  {
    \file_if_exist:nT { /System/Library/Fonts/Menlo.ttc }
      { \bool_gset_true:N \g__hit_platform_mac_bool }
  }
%    \end{macrocode}
%
% \changes{v3.2a}{0000/00/00}{\option{font} 的 \texttt{auto} 一档改成除开源字体集外
%   一律优先 Times New Roman。原先只有 \option{fontset} 取 \texttt{windows} 或
%   \texttt{mac} 时才用，\texttt{adobe} 与 \texttt{founder} 这两个 Windows 上的
%   商业字体集反而拿不到 Times New Roman}
% \option{font}\texttt{=auto} 的规则：\option{fontset} 写的是开源字体集
% （\texttt{fandol}、\texttt{ubuntu}）就配开源西文，其余一律优先
% Times New Roman，机器上没装再退回 TeX Gyre Termes。
%
% 反过来列开源集有两个好处。一是 \pkg{ctex} 的 \texttt{macnew}、\texttt{macold}
% 这些没列进来的值不会掉进错误的兜底；二是写 \texttt{fandol} 或 \texttt{ubuntu}
% 的人多半图的就是"到哪台机器都编出一样的东西"，给他配一个装没装看运气的
% Times New Roman 不合他的意思。回归测试的 42 个变体全用 \texttt{fandol}，
% 靠的也是这条。
%
% 平台探测在这条路上不再需要：不写 \option{fontset} 时同样先试
% Times New Roman，没有就退回，结果与原先按平台猜一致。
%    \begin{macrocode}
\str_set:Nx \l__hit_font_str { \hit@font }
\str_if_eq:VnT \l__hit_font_str { } { \str_set:Nn \l__hit_font_str { auto } }
\str_if_eq:VnT \l__hit_font_str { auto }
  {
    \str_case:enF { \hit@fontset }
      {
        { fandol } { \str_set:Nn \l__hit_font_str { termes } }
        { ubuntu } { \str_set:Nn \l__hit_font_str { termes } }
      }
      { \str_set:Nn \l__hit_font_str { times } }
  }
\str_if_eq:VnT \l__hit_font_str { times }
  {
    \fontspec_font_if_exist:nF { Times~New~Roman }
      { \str_set:Nn \l__hit_font_str { termes } }
  }
\cs_new_protected:Npn \__hit_set_font_times:
  {
    \setmainfont { Times~New~Roman }
    \setsansfont { Arial }
    \setmonofont { Courier~New } [ Scale = MatchLowercase ]
  }
\cs_new_protected:Npn \__hit_set_font_termes:
  {
    \setmainfont { TeXGyreTermesX }
      [
        UprightFont    = *-Regular ,
        BoldFont       = *-Bold ,
        ItalicFont     = *-Italic ,
        BoldItalicFont = *-BoldItalic ,
        Extension      = .otf ,
      ]
    \setsansfont { texgyreheros }
      [
        UprightFont    = *-regular ,
        BoldFont       = *-bold ,
        ItalicFont     = *-italic ,
        BoldItalicFont = *-bolditalic ,
        Extension      = .otf ,
      ]
    \setmonofont { texgyrecursor }
      [
        UprightFont    = *-regular ,
        BoldFont       = *-bold ,
        ItalicFont     = *-italic ,
        BoldItalicFont = *-bolditalic ,
        Extension      = .otf ,
      ]
  }
\cs_new_protected:Npn \__hit_set_font_none: { }
%    \end{macrocode}
%
% \changes{v3.2a}{0000/00/00}{\option{font} 的取值写错时报错。原先是裸的
%   \cs{use:c}，\cs{csname} 会把没定义的名字变成 \cs{relax}，于是拼错的人拿到
%   一份完全没设西文字体的文档，还看不到任何提示}
% 与 \option{cjk-font} 一致，先查函数在不在，不在就报错。
%    \begin{macrocode}
\msg_new:nnn { hithesis } { font-unknown }
  { 选项~font~的取值~'##1'~无效。 }
\cs_if_exist:cTF { __hit_set_font_ \l__hit_font_str : }
  { \use:c { __hit_set_font_ \l__hit_font_str : } }
  { \msg_error:nnx { hithesis } { font-unknown } { \l__hit_font_str } }
\cs_new_protected:Npn \__hit_set_cjk_font_windows:
  {
    \setCJKmainfont { SimSun } [ AutoFakeBold = true , ItalicFont = KaiTi ]
    \setCJKsansfont { SimHei } [ AutoFakeBold = true ]
    \setCJKmonofont { FangSong }
    \setCJKfamilyfont { zhsong } { SimSun } [ AutoFakeBold = true ]
    \setCJKfamilyfont { zhhei }  { SimHei } [ AutoFakeBold = true ]
    \setCJKfamilyfont { zhkai }  { KaiTi }
    \setCJKfamilyfont { zhfs }   { FangSong }
  }
\cs_new_protected:Npn \__hit_set_cjk_font_mac:
  {
    \defaultCJKfontfeatures { }
    \setCJKmainfont { Songti~SC }
      [
        UprightFont    = *~Light ,
        BoldFont       = *~Bold ,
        ItalicFont     = Kaiti~SC~Regular ,
        BoldItalicFont = Kaiti~SC~Bold ,
      ]
    \setCJKsansfont { Heiti~SC } [ UprightFont = *~Light , BoldFont = *~Medium ]
    \setCJKmonofont { STFangsong }
    \setCJKfamilyfont { zhsong } { Songti~SC }
      [ UprightFont = *~Light , BoldFont = *~Bold ]
    \setCJKfamilyfont { zhhei } { Heiti~SC }
      [ UprightFont = *~Light , BoldFont = *~Medium ]
    \setCJKfamilyfont { zhfs }  { STFangsong }
    \setCJKfamilyfont { zhkai } { Kaiti~SC }
      [ UprightFont = *~Regular , BoldFont = *~Bold ]
  }
\cs_new_protected:Npn \__hit_set_cjk_font_ubuntu:
  {
    \defaultCJKfontfeatures { }
    \setCJKmainfont { Noto~Serif~CJK~SC }
      [
        UprightFont    = *~Light ,
        BoldFont       = *~Bold ,
        ItalicFont     = FandolKai-Regular ,
        ItalicFeatures = { Extension = .otf } ,
        Script         = CJK ,
      ]
    \setCJKsansfont { Noto~Sans~CJK~SC } [ BoldFont = *~Medium , Script = CJK ]
    \setCJKmonofont { Noto~Sans~Mono~CJK~SC } [ Script = CJK ]
    \setCJKfamilyfont { zhsong } { Noto~Serif~CJK~SC }
      [ UprightFont = *~Light , BoldFont = *~Bold , Script = CJK ]
    \setCJKfamilyfont { zhhei } { Noto~Sans~CJK~SC }
      [ BoldFont = *~Medium , Script = CJK ]
    \setCJKfamilyfont { zhfs }  { FandolFang } [ Extension = .otf , UprightFont = *-Regular ]
    \setCJKfamilyfont { zhkai } { FandolKai }  [ Extension = .otf , UprightFont = *-Regular ]
  }
\cs_new_protected:Npn \__hit_set_cjk_font_fandol:
  {
    \defaultCJKfontfeatures { }
    \setCJKmainfont { FandolSong }
      [
        Extension      = .otf ,
        UprightFont    = *-Regular ,
        BoldFont       = *-Bold ,
        ItalicFont     = FandolKai-Regular ,
        ItalicFeatures = { Extension = .otf } ,
      ]
    \setCJKsansfont { FandolHei } [ Extension = .otf , UprightFont = *-Regular , BoldFont = *-Bold ]
    \setCJKmonofont { FandolFang } [ Extension = .otf , UprightFont = *-Regular ]
    \setCJKfamilyfont { zhsong } { FandolSong }
      [ Extension = .otf , UprightFont = *-Regular , BoldFont = *-Bold ]
    \setCJKfamilyfont { zhhei } { FandolHei }
      [ Extension = .otf , UprightFont = *-Regular , BoldFont = *-Bold ]
    \setCJKfamilyfont { zhfs }  { FandolFang } [ Extension = .otf , UprightFont = *-Regular ]
    \setCJKfamilyfont { zhkai } { FandolKai }  [ Extension = .otf , UprightFont = *-Regular ]
  }
%    \end{macrocode}
%
% Adobe 与方正照抄 \pkg{ctex} 的同名 fontset，\option{cjk-font=auto} 也认这两个值。
%    \begin{macrocode}
\cs_new_protected:Npn \__hit_set_cjk_font_adobe:
  {
    \defaultCJKfontfeatures { }
    \setCJKmainfont { AdobeSongStd-Light }
      [ BoldFont = AdobeHeitiStd-Regular , ItalicFont = AdobeKaitiStd-Regular ]
    \setCJKsansfont { AdobeHeitiStd-Regular }
    \setCJKmonofont { AdobeFangsongStd-Regular }
    \setCJKfamilyfont { zhsong } { AdobeSongStd-Light }
    \setCJKfamilyfont { zhhei }  { AdobeHeitiStd-Regular }
    \setCJKfamilyfont { zhfs }   { AdobeFangsongStd-Regular }
    \setCJKfamilyfont { zhkai }  { AdobeKaitiStd-Regular }
  }
\cs_new_protected:Npn \__hit_set_cjk_font_founder:
  {
    \defaultCJKfontfeatures { }
    \setCJKmainfont { FZShuSong-Z01 }
      [ BoldFont = FZXiaoBiaoSong-B05 , ItalicFont = FZKai-Z03 ]
    \setCJKsansfont { FZXiHeiI-Z08 } [ BoldFont = FZHei-B01 ]
    \setCJKmonofont { FZFangSong-Z02 }
    \setCJKfamilyfont { zhsong } { FZShuSong-Z01 } [ BoldFont = FZXiaoBiaoSong-B05 ]
    \setCJKfamilyfont { zhhei }  { FZHei-B01 }
    \setCJKfamilyfont { zhfs }   { FZFangSong-Z02 }
    \setCJKfamilyfont { zhkai }  { FZKai-Z03 }
  }
\cs_new_protected:Npn \__hit_set_cjk_font_none: { }
%    \end{macrocode}
%
% \changes{v3.2a}{0000/00/00}{中文字体可由模板自己设置，不再只能转给 ctex}
% \changes{v3.2a}{0000/00/00}{\option{cjk-font} 的 \texttt{noto} 一档改名
%   \texttt{ubuntu}，与 \pkg{ctex} 的同名 fontset 对齐。这个选项 v3.2a 才加，
%   还没发过版，不走废弃期}
% 中文字体。\texttt{windows} 用中易，\texttt{mac} 用华文，\texttt{ubuntu} 用思源，
% \texttt{fandol} 用 Fandol，前四组抄自 \pkg{ustcthesis}。取值一律取
% \pkg{ctex} 的 fontset 名，所以思源那一档叫 \texttt{ubuntu} 而不是 \texttt{noto}。
%    \begin{macrocode}
\str_set:Nx \l__hit_cjk_font_str { \hit@cjkfont }
\str_if_eq:VnF \l__hit_cjk_font_str { }
  {
    \str_if_eq:VnT \l__hit_cjk_font_str { auto }
      {
        \str_case:enF { \hit@fontset }
          {
            { windows } { \str_set:Nn \l__hit_cjk_font_str { windows } }
            { mac }     { \str_set:Nn \l__hit_cjk_font_str { mac } }
            { ubuntu }  { \str_set:Nn \l__hit_cjk_font_str { ubuntu } }
            { fandol }  { \str_set:Nn \l__hit_cjk_font_str { fandol } }
            { adobe }   { \str_set:Nn \l__hit_cjk_font_str { adobe } }
            { founder } { \str_set:Nn \l__hit_cjk_font_str { founder } }
          }
          {
            \sys_if_platform_windows:TF
              { \str_set:Nn \l__hit_cjk_font_str { windows } }
              {
                \bool_if:NTF \g__hit_platform_mac_bool
                  { \str_set:Nn \l__hit_cjk_font_str { mac } }
                  { \str_set:Nn \l__hit_cjk_font_str { fandol } }
              }
          }
      }
    \msg_new:nnn { hithesis } { cjk-font-unknown }
      { 选项~cjk-font~的取值~'##1'~无效。 }
    \cs_if_exist:cTF { __hit_set_cjk_font_ \l__hit_cjk_font_str : }
      { \use:c { __hit_set_cjk_font_ \l__hit_cjk_font_str : } }
      { \msg_error:nnx { hithesis } { cjk-font-unknown } { \l__hit_cjk_font_str } }
%    \end{macrocode}
%
% \pkg{ctex} 只在自己管字体时才定义 \cs{songti} 这几个命令。现在字体由模板设置，
% 传给 \pkg{ctex} 的是 \texttt{fontset=none}，这几个命令要自己补。
% \texttt{cjk-font=none} 时同样定义，此时 \texttt{zhsong} 等四个字族由用户自己
% \cs{setCJKfamilyfont}。
%    \begin{macrocode}
    \cs_set_protected:Npn \songti { \CJKfamily { zhsong } }
    \cs_set_protected:Npn \heiti { \CJKfamily { zhhei } }
    \cs_set_protected:Npn \fangsong { \CJKfamily { zhfs } }
    \cs_set_protected:Npn \kaishu { \CJKfamily { zhkai } }
  }% \str_if_eq:VnF \l__hit_cjk_font_str { } 到此为止
  }% \hit@setup@fonts 到此为止
%    \end{macrocode}
%
% WORD 中的字号对应该关系如下（1bp = 72.27/72 pt）:
% \begin{center}
% \begin{tabular}{llll}
% \toprule
% 初号 & 42bp & 14.82mm & 42.1575pt \\
% 小初 & 36bp & 12.70mm & 36.135 pt \\
% 一号 & 26bp & 9.17mm & 26.0975pt \\
% 小一 & 24bp & 8.47mm & 24.09pt \\
% 二号 & 22bp & 7.76mm & 22.0825pt \\
% 小二 & 18bp & 6.35mm & 18.0675pt \\
% 三号 & 16bp & 5.64mm & 16.06pt \\
% 小三 & 15bp & 5.29mm & 15.05625pt \\
% 四号 & 14bp & 4.94mm & 14.0525pt \\
% 小四 & 12bp & 4.23mm & 12.045pt \\
% 五号 & 10.5bp & 3.70mm & 10.59375pt \\
% 小五 & 9bp & 3.18mm & 9.03375pt \\
% 六号 & 7.5bp & 2.56mm & \\
% 小六 & 6.5bp & 2.29mm & \\
% 七号 & 5.5bp & 1.94mm & \\
% 八号 & 5bp & 1.76mm & \\\bottomrule
% \end{tabular}
% \end{center}
%
% \changes{v3.2a}{0000/00/00}{\cs{hit@def@fontsize} 改名 \cs{\_\_hit\_def\_fontsize:nn}。
%   它只在本文件里用，按这一层自己的规矩该是内部名}
% 根据习惯定义字号。用法：\cs{\_\_hit\_def\_fontsize:nn}\marg{字号名称}\marg{磅数}，
% 避免字号选择和行距的紧耦合。所有字号定义时为单倍行距，并提供选项指定行距倍数。
%
% 底下这十七个是中文字号的标准阶梯，模板自己只用到十一个，另外六个
% （\cs{chuhao}、\cs{xiaochu}、\cs{liuhao}、\cs{xiaoliu}、\cs{qihao}、\cs{bahao}）
% 是留给用户在正文里用的，别看着"没引用"就删。v3.2a 删掉的那三个
% （\cs{dachu}、\cs{banxiaosi}、\cs{dawu}）不在标准阶梯里，那才是该删的。
%    \begin{macrocode}
\cs_new_protected:Npn \__hit_def_fontsize:nn #1#2
  {
    \exp_args:Nc \NewDocumentCommand {#1} { O{1.3} }
      { \fontsize {#2} { ##1 \dimexpr #2 } \selectfont }
  }
%    \end{macrocode}
%
%
% \changes{v3.2a}{0000/00/00}{删掉 \cs{dachu}、\cs{banxiaosi}、\cs{dawu} 三个字号。
%   它们不在中文字号的标准阶梯里，手册的字号对照表也没列，模板自己一处没用过}
% \begin{macro}{\chuhao}
% \begin{macro}{\xiaochu}
% \begin{macro}{\yihao}
% \begin{macro}{\xiaoyi}
% \begin{macro}{\erhao}
% \begin{macro}{\xiaoer}
% \begin{macro}{\sanhao}
% \begin{macro}{\xsanhao}
% \changes{v3.1b}{2022/06/06}{增加 \cs{xsanhao} 为了深圳校区的原创性声明标题。}
% \begin{macro}{\xiaosan}
% \begin{macro}{\sihao}
% \begin{macro}{\xiaosi}
% \begin{macro}{\wuhao}
% \begin{macro}{\xiaowu}
% \begin{macro}{\liuhao}
% \begin{macro}{\xiaoliu}
% \begin{macro}{\qihao}
% \begin{macro}{\bahao}
% 一组字号定义。
%    \begin{macrocode}
\__hit_def_fontsize:nn{chuhao}{42bp}
\__hit_def_fontsize:nn{xiaochu}{36bp}
\__hit_def_fontsize:nn{yihao}{26bp}
\__hit_def_fontsize:nn{xiaoyi}{24bp}
\__hit_def_fontsize:nn{erhao}{22bp}
\__hit_def_fontsize:nn{xiaoer}{18bp}
\__hit_def_fontsize:nn{sanhao}{16bp}
\__hit_def_fontsize:nn{xsanhao}{15.9bp}
\__hit_def_fontsize:nn{xiaosan}{15bp}
\__hit_def_fontsize:nn{sihao}{14bp}
\__hit_def_fontsize:nn{xiaosi}{12bp}
\__hit_def_fontsize:nn{wuhao}{10.5bp}
\__hit_def_fontsize:nn{xiaowu}{9bp}
\__hit_def_fontsize:nn{liuhao}{7.5bp}
\__hit_def_fontsize:nn{xiaoliu}{6.5bp}
\__hit_def_fontsize:nn{qihao}{5.5bp}
\__hit_def_fontsize:nn{bahao}{5bp}
%    \end{macrocode}
%
% \cs{xsanhao} 是给深圳校区本科生的原创性声明标题用的。规范里那个标题是三号，
% 但 \cs{sanhao} 一行排不下，所以参考 \hitszthesis{} 定义了一个略小的字号。
% 只有 \file{hithesisbook.cls} 用得到，三个类都定义是因为多一条 \cs{fontsize}
% 不费什么，没必要为它单开一处特判。
%    \begin{macrocode}
%</shared-fonts>
%    \end{macrocode}
%
% \end{macro}
% \end{macro}
% \end{macro}
% \end{macro}
% \end{macro}
% \end{macro}
% \end{macro}
% \end{macro}
% \end{macro}
% \end{macro}
% \end{macro}
% \end{macro}
% \end{macro}
% \end{macro}
% \end{macro}
% \end{macro}
% \end{macro}
%
% \Finale
